% A UML sequence diagram of the issue-delegate-verify ceremony, drawn from the
% chapter's own Phase-3 pseudocode: Daemon.IssueToken (Alg. asymmetric),
% Agent.Delegate and Harbor.VerifyChain (Alg. delegation), with the capability
% sets of Def. offline-attn (["db:read","db:write"] narrowed to ["db:read"]).
%
% Keys, per signer (the 2026-11 correction; the earlier drawing verified every
% hop under pk_daemon, which contradicted Alg. delegation). card_0 is signed by
% the daemon and NAMES its subject A and A's public key pk_A; card_1 is signed
% by A and names B and pk_B. The Harbor verifies card_0 under pk_daemon and
% card_1 under the pk_A that the authenticated card_0 names -- exactly the
% pk_current walk of Harbor.VerifyChain. The ProVerif models v6/v7 likewise
% check each delegation under its signer's key (a_pk, b_pk), though they pass
% those keys in as process parameters instead of reading them from the parent
% card; the chapter says so beside the v6/v7 excerpt.
%
% Conventions in force, so that nothing in the drawing varies except what the
% protocol varies: one arrow style for every communication; activation boxes
% on lifelines for "this participant is working"; a reflex self-message for
% the verifier's own loop; arrowheads stopped 0.16cm short of the box they
% address; message labels wrapped to sit inside their own lane, so they need
% no opaque backing and cross no lifeline -- which is why `pd direct label`
% is used here with fill=none throughout.
%
% The readable fact is the shaded region: every message after m1 is inside it,
% and the Daemon's lifeline crosses it carrying no activation box at all.
% m1 now carries a five-line label (subject and subject key added), so m1 and
% the Daemon's activation box sit 0.30cm lower than before and the region's
% top edge 0.15cm lower; m2 (now five lines, with pk_B) sits 0.20cm lower so
% its label starts inside the region. The verifier's note now names both keys
% on five short lines, so its reflex, m4 and the region's foot sit 0.50cm lower. Re-checked
% on a standalone render at 110 dpi (Nov 2026): no label crosses the region
% edge or a lifeline, and the Daemon has no activation box inside the region.
% NO dash-bearing path exists in this figure, deliberately: an earlier version
% promised a dotted lifeline that rendered solid, and the fix was to stop
% encoding the fact in a dash at all.
% [internal; the chapter's own algorithms and the ProVerif models]
\begin{figure}[H]
\centering
\begin{tikzpicture}[pd figure,x=1cm,y=1cm,
  msg/.style={pd direct label,fill=none,align=center,anchor=south,inner ysep=3pt}]

  % ---- the offline region, drawn first so every mark sits on top of it ----
  % Encoded, not annotated: the reader sees a region with a drawn edge, and
  % that edge is itself the divider under message 1.
  \draw[pd neutral fill] (-5.85,-13.05) rectangle (5.00,-3.55);
  % The region's name runs up its own left gutter -- inside the region it
  % names, clear of every lifeline, and not mistakable for a note set loose
  % under the drawing.
  \node[pd note,rotate=90,anchor=south,align=center] at (-5.73,-8.30)
    {offline: the Daemon is not contacted again};

  % ---- participants ----
  \node[pd actor,text width=1.35cm,align=center] at (-4.72,0) {Daemon};
  \node[pd actor,text width=1.35cm,align=center] at (-1.77,0) {Agent A};
  \node[pd actor,text width=1.35cm,align=center] at ( 1.18,0) {Agent B};
  \node[pd actor,text width=1.35cm,align=center] at ( 4.13,0) {Harbor};

  % ---- lifelines: one style for all four. Who is idle is shown by the
  %      absence of an activation box, not by a second line style. ----
  \foreach \x in {-4.72,-1.77,1.18,4.13} \draw[pd hairline] (\x,-0.78) -- (\x,-12.85);

  % ---- activation boxes ----
  % The Daemon's is the only one above the region, and it has none below it.
  \draw[pd rule,fill=hhpaper] (-4.83,-2.65) rectangle (-4.61,-3.15);
  \draw[pd rule,fill=hhpaper] (-1.88,-2.90) rectangle (-1.66,-5.80);
  \draw[pd rule,fill=hhpaper] ( 1.07,-5.80) rectangle ( 1.29,-8.60);
  \draw[pd rule,fill=hhpaper] ( 4.02,-8.60) rectangle ( 4.24,-12.10);
  \draw[pd rule,fill=hhpaper] ( 1.07,-12.10) rectangle ( 1.29,-12.50);

  % ---- m1: the Daemon issues the root card, naming A and A's key ----
  \draw[pd arrow] (-4.61,-2.90) -- (-2.04,-2.90);
  \node[msg] at (-3.245,-2.78)
    {\texttt{m1\quad card\_0 =}\\\texttt{Sign(sk\_daemon,}\\\texttt{jti, A, pk\_A,}\\\texttt{\{db:read,}\\\texttt{db:write\})}};

  % ---- m2: agent A signs the child card with its own key, no round trip ----
  \draw[pd arrow] (-1.66,-5.80) -- (0.91,-5.80);
  \node[msg] at (-0.295,-5.68)
    {\texttt{m2\quad card\_1 =}\\\texttt{Sign(sk\_A,}\\\texttt{card\_0, B,}\\\texttt{pk\_B,}\\\texttt{\{db:read\})}};

  % ---- m3: agent B presents the whole chain ----
  \draw[pd arrow] (1.29,-8.60) -- (3.86,-8.60);
  \node[msg] at (2.655,-8.48)
    {\texttt{m3\quad chain =}\\\texttt{[card\_0,}\\\texttt{card\_1]}};

  % ---- the verifier's own loop: a reflex self-message on its activation box --
  \draw[pd arrow] (4.02,-9.40) -- (3.62,-9.40) -- (3.62,-10.40) -- (3.86,-10.40);
  \node[pd direct label,fill=none,align=right,anchor=east,text width=2.05cm]
    at (3.49,-9.90) {verify \texttt{card\_0}\\ under\\ \texttt{pk\_daemon},\\ \texttt{card\_1} under\\ \texttt{pk\_A}};

  % ---- m4: the verdict returns to B, never to the Daemon ----
  \draw[pd arrow] (4.02,-12.10) -- (1.45,-12.10);
  \node[msg] at (2.655,-11.98)
    {\texttt{m4\quad accept}\\\texttt{\{db:read\}}};
\end{tikzpicture}
\caption{The issue--delegate--verify ceremony as a sequence diagram. Every
message after \texttt{m1} falls inside the shaded region, where the Daemon is
not addressed. Each card names its subject and the subject's public key.
Agent~A signs \texttt{card\_1} with its own key, and the Harbor checks each
card under its signer's key: \texttt{card\_0} under \texttt{pk\_daemon}, and
\texttt{card\_1} under \texttt{pk\_A}, the key that \texttt{card\_0} names. It
never contacts the issuer. [internal]}
\label{fig:anchor-handshake-ladder}
\end{figure}
